\documentclass[11pt]{article}
\usepackage[utf8]{inputenc}
\usepackage{amsmath,amssymb}
\usepackage{hyperref}
\title{Robust Polygenic Risk Score Prediction under Distribution Shift}
\author{Anonymous}
\date{}
\begin{document}
\maketitle

\begin{abstract}
This paper studies Polygenic Risk Score Prediction. Grounded in a real, citable literature \cite{ref1,ref2,ref3,ref4,ref5,ref6,ref7,ref8},
we propose a focused method and report \textbf{illustrative} experiments.
All references below are real and verifiable (OpenAlex/arXiv); the reported
numbers are simulated to illustrate intended behaviour.
\end{abstract}

\section{Introduction}
Polygenic Risk Score Prediction is an active research area. This work targets a concrete gap identified from
the recent literature and contributes a method together with an evaluation protocol.

\section{Related Work (real literature)}
The following works, retrieved from public bibliographic APIs, frame our study:
\begin{itemize}
\item Khera et al. (2018) — "Genome-wide polygenic scores for common diseases identify individuals with risk equivalent to monogenic mutations", Nature Genetics (cited 3143).
\item Ge et al. (2019) — "Polygenic prediction via Bayesian regression and continuous shrinkage priors", Nature Communications (cited 2025).
\item Torkamani et al. (2018) — "The personal and clinical utility of polygenic risk scores", Nature Reviews Genetics (cited 1762).
\item Dudbridge (2013) — "Power and Predictive Accuracy of Polygenic Risk Scores", PLoS Genetics (cited 1658).
\item Michailidou et al. (2017) — "Association analysis identifies 65 new breast cancer risk loci", Nature (cited 1599).
\item Martin et al. (2017) — "Human Demographic History Impacts Genetic Risk Prediction across Diverse Populations", The American Journal of Human Genetics (cited 1583).
\end{itemize}

\section{Method}
We model distribution shift explicitly and optimise a worst-case objective over an uncertainty set of plausible test distributions. We instantiate this idea for Polygenic Risk Score Prediction and analyse its cost/quality trade-off.

\section{Experiments (Illustrative)}
\begin{center}
\begin{tabular}{lcc}
System & Primary Metric & Efficiency \\ \hline
Strong baseline & 0.812 & 1.00$\times$ \\
Ablation & 0.828 & 0.92$\times$ \\
\textbf{Ours} & \textbf{0.851} & \textbf{0.71$\times$}
\end{tabular}
\end{center}
\emph{Metrics simulated to illustrate the intended effect; not measured.}

\section{Conclusion}
We connected a real literature to a concrete method for Polygenic Risk Score Prediction. Future work will
validate the illustrative results with real experiments.

\begin{thebibliography}{9}
\bibitem{ref1} Amit V. Khera, Mark Chaffin, Krishna G. Aragam et al.. Genome-wide polygenic scores for common diseases identify individuals with risk equivalent to monogenic mutations. \emph{Nature Genetics}, 2018. DOI:10.1038/s41588-018-0183-z
\bibitem{ref2} Tian Ge, Chia‐Yen Chen, Yang Ni et al.. Polygenic prediction via Bayesian regression and continuous shrinkage priors. \emph{Nature Communications}, 2019. DOI:10.1038/s41467-019-09718-5
\bibitem{ref3} Ali Torkamani, Nathan E. Wineinger, Eric J. Topol. The personal and clinical utility of polygenic risk scores. \emph{Nature Reviews Genetics}, 2018. DOI:10.1038/s41576-018-0018-x
\bibitem{ref4} Frank Dudbridge. Power and Predictive Accuracy of Polygenic Risk Scores. \emph{PLoS Genetics}, 2013. DOI:10.1371/journal.pgen.1003348
\bibitem{ref5} Kyriaki Michailidou, Sara Lindström, Joe Dennis et al.. Association analysis identifies 65 new breast cancer risk loci. \emph{Nature}, 2017. DOI:10.1038/nature24284
\bibitem{ref6} Alicia R. Martin, Christopher R. Gignoux, Raymond K. Walters et al.. Human Demographic History Impacts Genetic Risk Prediction across Diverse Populations. \emph{The American Journal of Human Genetics}, 2017. DOI:10.1016/j.ajhg.2017.03.004
\bibitem{ref7} Nasim Mavaddat, Kyriaki Michailidou, Joe Dennis et al.. Polygenic Risk Scores for Prediction of Breast Cancer and Breast Cancer Subtypes. \emph{The American Journal of Human Genetics}, 2018. DOI:10.1016/j.ajhg.2018.11.002
\bibitem{ref8} Daifeng Wang, Shuang Liu, Jonathan Warrell et al.. Comprehensive functional genomic resource and integrative model for the human brain. \emph{Science}, 2018. DOI:10.1126/science.aat8464
\end{thebibliography}
\end{document}
